\chapter{Complexity — Layout}
%%% generated by the blueprint pipeline from TCSlib.Complexity.SpaceComplexity.Codec.Layout
\section{Overview}

Declarations of the Lean module \texttt{TCSlib.Complexity.SpaceComplexity.Codec.Layout}.

\section{Declarations}

\begin{lemma}[Bound on the binary-length power]
\lean{Complexity.Codec.two_pow_bits_le}
\label{Complexity.Codec.two_pow_bits_le}
\leanok
\difficulty{4}
For every natural number $L\ge1$, if $b$ is the length of the binary expansion of $L$,
then $2^b\le 2L$.
\end{lemma}

\begin{lemma}[Bound on the binary width]
\lean{Complexity.Codec.wP_le}
\label{Complexity.Codec.wP_le}
\leanok
\uses{Complexity.Codec.wP, Complexity.LogProg.length_bits_le_log}
\difficulty{3}
For every natural number $n$, the length of the binary expansion of $n+1$ is at most
$n+1$.
\end{lemma}

\begin{lemma}[Recovering a bit from a nested encoding]
\lean{Complexity.Codec.getD_w}
\label{Complexity.Codec.getD_w}
\leanok
\uses{Complexity.Codec.pe_getD_fst, Complexity.Codec.pe_getD_snd, Turing.pairEncode}
\difficulty{1}
Let $x$, $w$, and $r$ be finite strings of bits, and let $i$ be a natural number with
$i<|w|$. In the encoding of $x$ paired with the encoding of $w$ paired with $r$, the bit
at position $2|x|+2+2i$, taking the default value false if that position is out of
range, is equal to the bit at position $i$ of $w$, again taking false as the default.
\end{lemma}

\begin{lemma}[Configuration encoding bound]
\lean{Complexity.Codec.bnd_ge}
\label{Complexity.Codec.bnd_ge}
\leanok
\uses{Complexity.Codec.W, Complexity.Codec.bnd}
\difficulty{4}
Let $N$ be a finite nondeterministic Turing machine with Boolean tape symbols, let $y$
be a finite Boolean string, and let $s$ be a natural number with $s\le 2|y|$. Write $H$
for the number of pairs consisting of an optional state of $N$ and an element of a
three-element set, and set $W=H+3k+1$, where $k$ is the number of work tapes. If
$b(y)=(8W+8)(|y|+1)$, then $|y|+2W(2s+1)+2(|y|+1)+4\le b(y)$.
\end{lemma}

\begin{lemma}[Coordinate of a conditional index]
\lean{Complexity.Codec.Regs.val_ite}
\label{Complexity.Codec.Regs.val_ite}
\leanok
\uses{Complexity.Codec.Regs, Complexity.Codec.Regs.val}
\difficulty{1}
For every register record $R$, proposition $c$ with a decision procedure, and indices
$r,r' \in\mathbb Z$, the coordinate value at the index selected by the condition is
the corresponding conditional value: $R(r)$ if $c$ holds, and $R(r')$ otherwise.
\end{lemma}

\begin{definition}[Register layout for a scan]
\lean{Complexity.Codec.layR}
\label{Complexity.Codec.layR}
\leanok
\uses{Complexity.Codec.Regs, Complexity.Codec.W, Complexity.Codec.wP}
Let $N$ be a finite nondeterministic Turing machine, let $R$ be a register record, and
let $n,s,pos$ be natural numbers. The scan layout is the register record obtained from
$R$ by setting registers $N1$ and $Wp$ to $n+1$ and $w_P(n)$, respectively; setting $E$
and $Q1$ to $2n+2+4W(N)s+2W(N)$, $F$ to that quantity plus $2w_P(n)$, and $U$ to that
quantity plus $2w_P(n)+4$; and setting $Pos$ to $pos$. Every other register retains its
value from $R$.
\end{definition}

\begin{theorem}[Layout run and rejection]
\lean{Complexity.Codec.layout_run}
\label{Complexity.Codec.layout_run}
\leanok
\uses{Complexity.Codec.CA, Complexity.Codec.Reach, Complexity.Codec.Reach.cast, Complexity.Codec.Reach.single', Complexity.Codec.Reach.trans, Complexity.Codec.Regs, Complexity.Codec.Regs.Bd, Complexity.Codec.Regs.Bd.val_le, Complexity.Codec.Regs.val, Complexity.Codec.Regs.val_E, Complexity.Codec.Regs.val_ite, Complexity.Codec.W, Complexity.Codec.bitsVal_lt_two_pow, Complexity.Codec.bnd, Complexity.Codec.bnd_ge, Complexity.Codec.cpDst, Complexity.Codec.cpSrc, Complexity.Codec.cp_run, Complexity.Codec.dcAcc, Complexity.Codec.dcBase, Complexity.Codec.dcPtr, Complexity.Codec.dcStride, Complexity.Codec.dc_run, Complexity.Codec.gd_of, Complexity.Codec.getD_w, Complexity.Codec.halt_rej, Complexity.Codec.hc_run, Complexity.Codec.icCount, Complexity.Codec.icL, Complexity.Codec.icReg, Complexity.Codec.ic_run, Complexity.Codec.layR, Complexity.Codec.orc, Complexity.Codec.orc_bit, Complexity.Codec.rB, Complexity.Codec.rP, Complexity.Codec.rU, Complexity.Codec.scDst, Complexity.Codec.scFactor, Complexity.Codec.scSrc, Complexity.Codec.sc_run, Complexity.Codec.two_pow_bits_le, Complexity.Codec.wP, Complexity.Codec.wP_le, Complexity.LogProg.AHalt, Complexity.LogProg.GuardLe, Complexity.LogProg.PreS, Complexity.LogProg.astep, Complexity.LogProg.bitsVal, Complexity.LogProg.vword_whole₁, Turing.length_pairEncode, Turing.pairEncode}
\difficulty{7}
Let $N$ be a finite nondeterministic Turing machine with Boolean input, let the four bit
strings $y,x,w,\mathrm{rest}$ satisfy $y=\langle x,\langle
w,\mathrm{rest}\rangle\rangle$ under the self-delimiting pair encoding, and let $R$ be a
vector of 21 natural-number registers. Suppose every register of $R$ is at most
$(8W(N)+8)(|y|+1)$, and that its registers $L,A,S,N_1,W_p,E,F,T,U,\mathrm{Pos}$ have
values $|y|,2|x|+2,s,0,0,0,0,0,0,0$, respectively, for some $s\le 2|y|$. Also suppose
the value of register $B$ in step mode, or register $P$ otherwise, is $2|x|+2|w|+4$.
Then either $w$ has length $W(N)(2s+1)+w_P(|x|)+1$ and its bit at position
$W(N)(2s+1)+w_P(|x|)$ is $1$, in which case the constructed register machine reaches its
comparison phase from the copy phase while all registers stay within the stated bound,
with its final register values given by the specified layout; or that length-and-bit
condition fails and the machine halts with answer $0$, again with all registers staying
within the bound.
\end{theorem}

\begin{theorem}[Comparison phase outcome]
\lean{Complexity.Codec.le_phase}
\label{Complexity.Codec.le_phase}
\leanok
\uses{Complexity.Codec.CA, Complexity.Codec.Reach, Complexity.Codec.Reach.single', Complexity.Codec.Reach.trans, Complexity.Codec.Regs, Complexity.Codec.Regs.Bd, Complexity.Codec.Regs.Bd.val_le, Complexity.Codec.Regs.upd_T, Complexity.Codec.Regs.val, Complexity.Codec.Regs.val_N1, Complexity.Codec.Regs.val_Pos, Complexity.Codec.bnd, Complexity.Codec.gd_of, Complexity.Codec.halt_rej, Complexity.Codec.le_run, Complexity.Codec.orc, Complexity.LogProg.AHalt, Complexity.LogProg.GuardLe, Complexity.LogProg.PreS, Complexity.LogProg.astep}
\difficulty{5}
Let $N$ be a finite nondeterministic Turing machine with Boolean input, let $y$ be a bit
string, and let $R$ be a vector of 21 natural-number registers whose values are all at
most $(8W(N)+8)(|y|+1)$. Suppose register $T$ has value $0$. If the value of register
$\mathrm{Pos}$ is at most that of register $N_1$, the constructed register machine
reaches its designated next phase—$E2$ in step mode and the clearing phase
otherwise—without changing the register valuation, while satisfying the bound and the
machine's precondition throughout. If $\mathrm{Pos}>N_1$, it instead halts with answer
$0$, also satisfying those conditions throughout its run.
\end{theorem}

\begin{theorem}[Execution of the configuration scan]
\lean{Complexity.Codec.vscan_run}
\label{Complexity.Codec.vscan_run}
\leanok
\uses{Complexity.Codec.CA, Complexity.Codec.Reach, Complexity.Codec.Reach.cast, Complexity.Codec.Reach.single', Complexity.Codec.Reach.trans, Complexity.Codec.Regs, Complexity.Codec.Regs.Bd, Complexity.Codec.Regs.Bd.val_le, Complexity.Codec.Regs.val, Complexity.Codec.VAd, Complexity.Codec.W, Complexity.Codec.bnd, Complexity.Codec.cpDst, Complexity.Codec.cpSrc, Complexity.Codec.cp_run, Complexity.Codec.gd_of, Complexity.Codec.getD_w, Complexity.Codec.vAd0, Complexity.Codec.vScan, Complexity.Codec.v_run, Complexity.LogProg.PreS, Complexity.LogProg.astep, Turing.pairEncode}
\difficulty{5}
Let $N$ be a finite nondeterministic Turing machine with Boolean input and work-tape
symbols, and fix any codec mode. Let $y,x,w,r$ be finite bit strings with $y$ equal to
the self-delimiting pair encoding of $x$ and the pair encoding of $w$ with $r$. Let $R$
be a register assignment whose 21 register values are all at most $B=(8W_N+8)(|y|+1)$,
and suppose its registers $A$ and $E$ have values $2|x|+2$ and $2|x|+2+4W_Ns+2W_N$,
respectively. If $W_N(2s+1)\le |w|$, then there is an accumulator state whose scan
component equals the result of scanning the first $W_N(2s+1)$ bits of $w$, and the codec
program in the chosen mode can reach its completion state from the register state
obtained by clearing $Q_1$, while keeping every register at most $B$ throughout the run;
at completion, registers $Q_1$ and $Q_2$ both equal the original value of $E$.
\end{theorem}
